\section{Experimental setup}\label{sec:setup}

Every condition shares the student initialization, cold-start data, prompts, seeds, sampling
settings and evaluation harness; in the main grid only the teacher's precision changes, and no
hyperparameter is re-tuned per teacher rung.

\subsection{Models}\label{sec:setup-models}
All models are from the Qwen3 family \citep{qwen3}, so teacher and student share a tokenizer,
which token-level OPD requires. The main teacher is Qwen3-30B-A3B in thinking mode (a
mixture-of-experts model with 30.5B total and 3.3B active parameters). The dense Qwen3-8B serves
as a second teacher for the equal-memory comparison and separates mixture-of-experts from dense
quantization effects. The students are Qwen3-1.7B (S1, main student) and Qwen3-0.6B (S2, cheap
student for seeds and replication), starting from the post-trained checkpoints
\todo{confirm post-trained vs.\ Base at the 1.4 checkpoint}. Every model is pinned to a Hub
commit. Before use, a hard check compares vocabulary, merges, added and special tokens, chat
template and the tokenization pipeline of every teacher--student pair and refuses to run on any
difference (S1 vs.\ S2 verified; teacher pairs \TBD).

\subsection{Teacher quantization ladder}\label{sec:setup-ladder}
Table~\ref{tab:ladder} lists the rungs. All quantized rungs are calibrated on the same
\nCalib{} prompts, drawn from the training pool but never trained on, each paired with a
reasoning trace written by the BF16 teacher in chat format and cut at 2{,}048 tokens. We record
each rung's memory, load time, scoring throughput and its own MATH-500 accuracy; the latter is
the ``teacher accuracy'' predictor in H2.

% generated by scripts/make_paper_tables.py from configs/models, configs/quant; do not edit
\begin{table}[t]
  \centering\small
  \caption{Teacher quantization ladder. Memory and MATH-500 accuracy are measured per rung in Phase 4.}\label{tab:ladder}
  \begin{tabular}{llllllrr}
    \toprule
    Rung & Model & Format & Bits & Group & Tool & Memory (GB) & MATH-500 \\
    \midrule
    T16 & Qwen3-30B-A3B & BF16 & 16 & -- & -- & 61.1 & 96.6 \\
    T8 & Qwen3-30B-A3B & W8A8-FP8 & 8 & -- & llmcompressor & 31.2 & 94.6 \\
    T4-fp & Qwen3-30B-A3B & NVFP4 & 4 & -- & llmcompressor & 18.1 & 93.6 \\
    T4-int & Qwen3-30B-A3B & W4A16 & 4 & 128 & llmcompressor & 16.7 & 94.2 \\
    T3 & Qwen3-30B-A3B & W3A16 & 3 & 128 & llmcompressor & 13.0 & 93.2 \\
    T2 & Qwen3-30B-A3B & Q2\_K & 2 & -- & llama.cpp & \TBD{} & \TBD{} \\
    D16 & Qwen3-8B & BF16 & 16 & -- & -- & \TBD{} & \TBD{} \\
    D8 & Qwen3-8B & W8A8-FP8 & 8 & -- & llmcompressor & \TBD{} & \TBD{} \\
    D4 & Qwen3-8B & W4A16 & 4 & 128 & llmcompressor & \TBD{} & \TBD{} \\
    \bottomrule
  \end{tabular}\par\smallskip\footnotesize{T2 is an optional stress rung. Memory and accuracy will be read from \texttt{reports/rungs.md}.}
\end{table}


\subsection{Data}\label{sec:setup-data}
\paragraph{Training prompts.} The training pool is DeepScaleR-Preview \citep{deepscaler2025}
(\nPoolRaw{} problems compiled from AIME 1984--2023, AMC before 2023, Omni-MATH and STILL).
We keep a problem only if math-verify \citep{mathverify} parses its answer to a single real
number with no free variables (integers, rationals and closed forms such as $\sqrt2-1$) and the
gold answer, written as \verb|\boxed{...}|, verifies against itself. We drop exact duplicates
(case- and whitespace-folded); when identical problems carry disagreeing answers we drop all
copies, since at least one label is wrong. Near-duplicates are removed with MinHash on word
5-grams (128 permutations, Jaccard $\ge 0.8$). This leaves \nPoolFiltered{} problems
(Table~\ref{tab:data}).

\paragraph{Contamination.} We drop every training prompt that shares a 13-gram (over lower-cased
alphanumeric word tokens, so markup and spacing cannot hide an overlap) with any evaluation item:
\nDecontDropped{} prompts, leaving \nPoolClean{}. Most matches are shared contest boilerplate
(``\dots where $m$ and $n$ are relatively prime positive integers. Find $m+n$''), which we drop
anyway; genuine near-copies (sharing at least half of an item's 13-grams) occur only for
MATH-500, where \nMathNearCopy{} evaluation problems appear in the training source.
\todo{GPQA-Diamond not yet checked: gated, pending access.}

\paragraph{Splits.} From the clean pool, with a fixed seed, we take \nTrain{} training prompts,
a cold-start subset of \nColdstart{} of them, \nCalib{} calibration prompts and \nProbe{} probe
prompts for the fidelity study; calibration and probe are disjoint from training and from each
other. A manifest records the hash of every file and of every split's id list.

% generated by scripts/make_paper_tables.py from data/pool/pool_stats.json, data/contamination/contamination_stats.json, data/manifest.json; do not edit
\begin{table}[t]
  \centering\small
  \caption{Training data pipeline (source: \texttt{agentica-org/DeepScaleR-Preview-Dataset}).}\label{tab:data}
  \begin{tabular}{lr}
    \toprule
    Stage & Problems \\
    \midrule
    Raw problems & 40{,}315 \\
    after normalize & 40{,}309 \\
    after answer\_type & 38{,}562 \\
    after self\_check & 38{,}560 \\
    after dedup\_exact & 37{,}199 \\
    after dedup\_near & 37{,}080 \\
    after 13-gram decontamination & 36{,}665 \\
    split: train & 16{,}000 \\
    split: coldstart & 5{,}000 \\
    split: calibration & 512 \\
    split: probe & 500 \\
    \bottomrule
  \end{tabular}
\end{table}


\paragraph{Evaluation.} Table~\ref{tab:evals} lists the benchmarks: GSM8K \citep{cobbe2021gsm8k},
MATH-500 \citep{hendrycks2021math,lightman2023verify}, AMC 2023, AIME 2024--2026 and HMMT
February 2026 \citep[the 2025/2026 sets from][]{matharena2025}, and GPQA-Diamond. AIME 2026 and HMMT February
2026 were released after the models and serve as a contamination-clean check; we report them
separately and say so if conclusions differ from the older sets. GPQA-Diamond \citep{rein2024gpqa}
is an out-of-domain check. Decoding uses thinking mode, temperature 0.6, top-p 0.95 and at most
16{,}384 tokens; we report pass@1 averaged over the samples per problem.

% generated by scripts/make_paper_tables.py from configs/eval/default.yaml, data/eval/eval_stats.json, data/contamination/contamination_stats.json; do not edit
\begin{table}[t]
  \centering\small
  \caption{Evaluation suites, samples per problem, and training prompts removed for 13-gram overlap with each suite.}\label{tab:evals}
  \begin{tabular}{lrrlr}
    \toprule
    Benchmark & Items & Samples & Source & Train dropped \\
    \midrule
    GSM8K & 1{,}319 & 1 & openai/gsm8k & 0 \\
    MATH-500 & 500 & 4 & HuggingFaceH4/MATH-500 & 241 \\
    AMC 2023 & 40 & 8 & math-ai/amc23 & 42 \\
    AIME 2024 & 30 & 16 & HuggingFaceH4/aime\_2024 & 228 \\
    AIME 2025 & 30 & 16 & MathArena/aime\_2025 & 202 \\
    AIME 2026 & 30 & 16 & MathArena/aime\_2026 & 209 \\
    HMMT Feb 2026 & 33 & 16 & MathArena/hmmt\_feb\_2026 & 1 \\
    GPQA-Diamond & \TBD{} & 4 & Idavidrein/gpqa & \TBD{} \\
    \bottomrule
  \end{tabular}
\end{table}


\paragraph{Answer checking.} Only the text after the last \verb|</think>| is graded; a response
whose thinking never closes (e.g.\ truncated) has no answer and counts as incorrect, and boxed
expressions inside the thinking block are ignored. For math, the last \verb|\boxed{}| is
compared with the gold by math-verify (falling back to math-verify's own extraction if nothing is
boxed). For multiple choice, the answer is the last option letter found by a fixed pattern
priority, starting with the requested \texttt{"answer": "C"} format. Every gold answer of the
built evaluation sets, written as a boxed final answer, is graded correct, and changing an
integer gold by one is always graded wrong.

\paragraph{Sanity check.} Table~\ref{tab:sanity} evaluates the official students with this
harness. \todo{Compare with the Qwen3 technical report (cite page); investigate any gap
$>5$ points before relying on the harness.}

% generated by scripts/make_paper_tables.py from runs/eval/S{1,2}__gsm8k+math500__n1__t16384__s0; do not edit
\begin{table}[t]
  \centering\small
  \caption{Sanity evaluation of the official students (thinking mode, 1 sample, 16{,}384 tokens): pass@1 and truncation rate (\%).}\label{tab:sanity}
  \begin{tabular}{lrrrrrr}
    \toprule
    Student & GSM8K & trunc. & MATH-500 & trunc. & Rep. GSM8K & Rep. MATH-500 \\
    \midrule
    S2 & 80.6 & 0.5 & 72.4 & 7.6 & n/a & 77.6 \\
    S1 & 90.1 & 0.2 & 86.0 & 7.8 & n/a & 93.4 \\
    \bottomrule
  \end{tabular}\par\smallskip\footnotesize{Rep.: Qwen3 technical report \citep{qwen3}, Table 19 (thinking mode, top-$k$ 20, up to 32{,}768 tokens); GSM8K is reported there only for base models.}
\end{table}


\subsection{Distillation procedure}\label{sec:setup-method}
\paragraph{Cold start.} Each student is first fine-tuned for one epoch on the \nColdstart{}
cold-start prompts answered once by the BF16 teacher in thinking mode, keeping only verified
correct traces. This set is generated once and shared by every condition.

\paragraph{OPD rounds.} One round (i) samples \nPrompts{} prompts $\times$ \nRollouts{}
rollouts from the student (temperature 1, at most \rolloutMaxTok{} tokens); (ii) scores every
sampled token with the teacher at bit-width $b$ through vLLM \citep{kwon2023vllm}
\texttt{prompt\_logprobs}, logging the top-\teacherTopK{} alternatives; (iii) forms
$r_t$ (Eq.~\ref{eq:reward}); (iv) takes one optimizer step on Eq.~\ref{eq:loss}; and (v) syncs
the weights into the rollout engine. The loss is averaged over all response tokens of the
round's \nPrompts{}$\times$\nRollouts{} rollouts,
\begin{equation}
  \hat{\mathcal{L}}(\theta) = -\frac{1}{N}\sum_{i}\sum_{t\in R_i} \sg(r_{i,t})\,
  \log\piS(y_{i,t}\mid x_i, y_{i,<t}),
  \label{eq:loss-batch}
\end{equation}
where $R_i$ are the response positions of rollout $i$ and $N=\sum_i|R_i|$; prompt and padding
positions carry no loss. When the round is split into micro-batches for gradient accumulation,
each micro-batch is normalized by the same $N$, so the accumulated gradient equals the
full-batch gradient. Unit tests check that the expected gradient of Eq.~\ref{eq:loss-batch}
equals $\nabla_\theta\mathrm{KL}(\piS\,\|\,\piT)$ on toy distributions. Runs last \nRounds{} rounds with a checkpoint (and a
MATH-500 evaluation) every 20 rounds. Table~\ref{tab:hparams} lists all settings; the learning
rate is chosen once, with the BF16 teacher, and then frozen for every rung.

\paragraph{Exact token handling.} Every rollout is stored as the exact prompt and response token
ids; the teacher scores exactly \texttt{prompt\_ids + response\_ids} and the scoring backends
assert this. Sampled sequences are not always the canonical tokenization of their own text, so
decoding and re-tokenizing would silently shift every per-token reward. The vLLM scoring backend
agrees with a plain transformers forward pass to a mean absolute difference of
\agreeMeanAbs{} nats per token (Appendix, Table~\ref{tab:agreement}).

\paragraph{Run modes.} Teacher and student are either co-resident in memory or loaded in
alternation each round; the choice is made from measured timings in a pilot run (E0).
\TBD{} \todo{mode chosen at the E0 gate, with the measured per-stage timings}

% generated by scripts/make_paper_tables.py from configs/train/{opd,sft}.yaml, configs/eval/default.yaml; do not edit
\begin{table}[t]
  \centering\small
  \caption{Training and evaluation settings, identical for every teacher rung.}\label{tab:hparams}
  \begin{tabular}{lp{5.2cm}p{4.2cm}}
    \toprule
    Setting & OPD & Cold-start SFT \\
    \midrule
    Optimizer & AdamW ($\beta$=0.9, 0.999; $\epsilon$=$10^{-8}$) & AdamW ($\beta$=0.9, 0.999; $\epsilon$=$10^{-8}$) \\
    Precision & BF16 mixed (FP32 weights and AdamW states) & BF16 mixed (FP32 weights and AdamW states) \\
    Weight decay / grad clip & 0 / 1 & 0 / 1 \\
    Learning rate & sweep 5e-7, 1e-6, 2e-6 (T16), then frozen & $10^{-5}$ \\
    Schedule & linear warmup 5 rounds, constant & warmup 3\%, cosine to 10\% \\
    Batch & 64 prompts $\times$ 4 rollouts, 1 step & 32 sequences \\
    Max length & 4{,}096 response tokens & 8{,}192 tokens (longer dropped) \\
    Rounds / epochs & 100 rounds & 1 epoch \\
    Rollout sampling & T=1, top-p=1 & -- \\
    Teacher top-k logged & 20 & -- \\
    Evaluation sampling & T=0.6, top-p=0.95, max 16{,}384 tokens & same \\
    \bottomrule
  \end{tabular}
\end{table}


\subsection{Baselines}
\emph{Cold start only} (evaluation of the starting checkpoints); \emph{off-policy SFT} on BF16
teacher traces, compute-matched to an OPD run; and \emph{RLVR} with GRPO \citep{shao2024grpo}
(verifier reward only, group size 4, wall-clock-matched). The RLVR baseline answers the practical
question of whether a cheap teacher is still better than no teacher.

\subsection{Statistics}\label{sec:stats}
\paragraph{Intervals.} For each condition we report the mean over seeds with a Student-$t$ 95\%
interval (two seeds, so these intervals are wide by design) and a percentile bootstrap interval
over problems (10{,}000 resamples), where pass@1 is the mean over problems of the per-problem
fraction of correct samples. Differences between two conditions on the same benchmark use a
paired bootstrap that resamples problems jointly.

\paragraph{H1 (equivalence).} For each quantized rung we compute, per MATH-500 problem, the
difference in pass rate (pooled over seeds and samples) between the rung's student and the
T16 student, and run two one-sided $t$-tests against a margin of $\pm1.0$ point, fixed before
E2. A rung is declared equivalent when both one-sided tests reject at $\alpha=0.05$, i.e.\ when
the 90\% interval of the mean difference lies inside $(-1, 1)$ points.
\todo{confirm the paired, per-problem formulation; a per-seed Welch TOST is also implemented}

\paragraph{H2 (predictor comparison).} Across conditions (rung $\times$ student) we compute the
Spearman correlation of the final student MATH-500 score with (a) the rung's fidelity score
(Section~\ref{sec:fidelity}) and (b) the rung's own MATH-500 accuracy, and bootstrap the
difference $|\rho_a|-|\rho_b|$ over conditions. With at most ten conditions this comparison is
low-powered, and we say so when reporting it.

\paragraph{Multiple comparisons.} The comparisons are fixed in advance---each quantized rung
against T16, and the two equal-memory teachers of E4---and corrected with Holm's step-down
procedure. TOST and Holm agree with the \texttt{statsmodels} implementations in our tests.

\subsection{Hardware}
All experiments run on one NVIDIA DGX Spark (Table~\ref{tab:hardware}); its 128\,GB of memory is
shared between CPU and GPU, and compute and memory bandwidth, not capacity, are the binding
limits.

% generated by scripts/make_paper_tables.py from reports/env.container.json; do not edit
\begin{table}[t]
  \centering\small
  \caption{Hardware and software (one DGX Spark).}\label{tab:hardware}
  \begin{tabular}{ll}
    \toprule
    Item & Value \\
    \midrule
    GPU & NVIDIA GB10, sm\_121 \\
    Memory (unified, GB) & 119.6 \\
    CPU & aarch64, 20 cores \\
    CUDA (container) & 13.4 \\
    PyTorch & 2.15.0a0+875d815502.nv26.10.72655024 \\
    vLLM & 0.31.0+88641915.nv26.10.75644909 \\
    Transformers & 5.17.0 \\
    \bottomrule
  \end{tabular}
\end{table}

